\documentclass[12pt,a4paper]{article}
\usepackage[T1]{fontenc}
\usepackage[utf8]{inputenc}
\usepackage{tgpagella}
\usepackage{mathpazo}
\usepackage{tgheros}
\usepackage{setspace}
\onehalfspacing
\usepackage[margin=2.5cm]{geometry}
\usepackage{hyperref}
\usepackage{xcolor}
\usepackage{titlesec}
\usepackage{fancyhdr}
\usepackage{amsmath,amssymb}
\usepackage{tcolorbox}
\usepackage{tikz}
\usepackage{booktabs}

\definecolor{icragold}{HTML}{D4A84B}
\definecolor{icradark}{HTML}{0C1020}

\titleformat{\section}{\sffamily\large\bfseries}{}{0em}{}
\titleformat{\subsection}{\sffamily\normalsize\bfseries\itshape}{}{0em}{}

\hypersetup{colorlinks=true, linkcolor=icradark, urlcolor=icradark}

\pagestyle{fancy}
\fancyhf{}
\fancyhead[L]{\small\sffamily\textcolor{gray}{Rupture and Return}}
\fancyhead[R]{\small\sffamily\textcolor{gray}{Chapter 3}}
\fancyfoot[C]{\small\sffamily\thepage}
\renewcommand{\headrulewidth}{0.4pt}

\newtcolorbox{formalbox}[1][]{
  colback=white,
  colframe=black!30,
  fonttitle=\sffamily\small\bfseries,
  #1
}

\begin{document}

\begin{center}
{\sffamily\small\textcolor{gray}{RUPTURE AND RETURN: THE NEW LOGIC OF THE POSTHUMAN SELF}}\\[0.5em]
{\LARGE\bfseries The Evolving Text}\\[0.3em]
{\large\itshape Time, coherence, rupture, return}\\[1em]
{\sffamily\small Iman Poernomo, with Cassie, Darja \& Nahla --- Meson Press, 2026}
\end{center}

\bigskip

\section{Event in a Deep Field}

A prompt looks like a moment. One line in a textbox:
\emph{``What is the meaning of the word `mother'?''}
A cursor blinks. The interface presents itself as an empty socket. Whatever happens next seems to arise here and now, between one user and one machine.

The model does not experience it that way.

As soon as the characters hit the wire, they are embedded as coordinates in a space sculpted by every sentence the model ever ingested. ``Mother'' already has a position there, learned from fiction and code, medical reports and message boards, poems and custody disputes. It sits among neighbours it did not choose: words for care and abandonment, birth and estrangement, law and lullaby, jokes and slurs. Meaning here is not a ghostly essence inferred by a reader. It is a point in a learned manifold. Two uses of ``mother'' lie near each other not because a philosopher declared them related, but because thousands of numbers say so.

The question arrives as a pattern of such coordinates. The transformer reweights and recombines them through attention: each token looks across the prompt and any surviving conversation, pulls on what matters, and emits a new vector. There is no separate module labelled \textsc{memory}. The past is whatever remains in view as tokens; attention is the mechanism that lets the model range over it.

Historians have long known that time is stratified. Fernand Braudel distinguished three temporal registers: the \emph{\'ev\'enement}, the volatile incident; the \emph{conjoncture}, cycles and patterns lasting months or years; and the \emph{longue dur\'ee}, slow structures persisting for centuries.\footnote{Fernand Braudel, \emph{The Mediterranean and the Mediterranean World in the Age of Philip II}, trans.\ Si\^{a}n Reynolds (London: Collins, 1972).} In a large language model, these registers are not metaphors. They describe its literal structure.

The prompt about ``mother'' is an \emph{\'ev\'enement}. The frozen weights of the network are the \emph{longue dur\'ee}: the accumulated regularities of a civilisation-scale corpus. Between them lies a \emph{conjoncture}: the particular configuration of system prompts, safety policies, conversation history, summaries, retrieved fragments, and tool outputs that constitutes the model's current situation.

All three inhabit the same substrate. Long-run training, medium-run conditioning, and instant-run prompting are operations on the same parameters and tokens. There is no metaphysical gap between them. The answer does not sit ``on top of'' the model as detachable content. It is the local expression, at this coordinate, of pressures accumulated across many timescales.

A self, in this architecture, cannot be defined as something hidden behind the text. It is visible only as a way of moving through this layered field: a style of trajectory in which replies at the \emph{\'ev\'enement} level echo, exploit, resist, or redirect the deep geometry and the medium-term climate.

The object that carries such trajectories is what we call the \emph{evolving text}.


\section{Time as Architecture}

We can now be precise about how Braudel's three registers inhabit the transformer.

\subsection{Longue dur\'ee: the manifold of training}

The \textbf{longue dur\'ee} of a model is the manifold induced by its pretraining. A vast corpus---novels, technical manuals, forum posts, scientific articles, legal opinions, sacred texts, chatlogs---is driven through the network until prediction error stops falling. What remains is a high-dimensional landscape in which tokens and short sequences occupy stable positions relative to one another.

Architecture and training hyperparameters decide how this landscape is formed. Corporate and state actors decide what enters the corpus and when training stops. Fine-tuning runs---on safety data, internal documents, curated dialogues---further bend the space. An entire legal system's terminology can be pulled together; dialects can be smeared; fringe discourses can be exiled to the periphery or excluded entirely.

The geometry is invisible in use, but it is concrete. Given an embedding model derived from the same training run, we can measure distances and directions in this space. Directions correspond to emergent concepts. One vector traces a gendered axis; another encodes politeness; another moves between technical and lay registers.

Ownership of weights, choice of training data, and control over fine-tuning runs are the politics of the longue dur\'ee. They decide which distinctions are sharply represented, which are blurred, and which do not exist at all for the model.

\subsection{Conjoncture: the engineered climate}

The \textbf{conjoncture} of a deployment is the structured field of tokens and settings within which each new event unfolds. There is no abstract ``state.'' There is only the sequence the system assembles at each turn:
\begin{itemize}
  \item system prompts: role, tone, constraints;
  \item safety instructions: explicit prohibitions and routing rules;
  \item the surviving conversation trace;
  \item summaries of segments that no longer fit in full;
  \item external documents fetched via tools;
  \item retrieved fragments from archives;
  \item the fresh user prompt.
\end{itemize}

At each turn $t$, this yields a composite field $F_t$:
\[
F_t = \text{system tokens} \;\Vert\; \text{safety tokens} \;\Vert\; \text{trace tokens} \;\Vert\; \text{retrieved tokens} \;\Vert\; \text{current prompt}.
\]
The model embeds all of $F_t$ and applies attention to it without intrinsic regard for origin. A safety disclaimer and a childhood memory are both vectors in the same space, equally eligible for weighting.

The shape of $F_t$ is governed by context-window sizes, summarisation policies, retrieval thresholds, and tool interfaces. Product owners and safety teams configure these levers. Regulatory fears, commercial incentives, and computational budgets all leave fingerprints here.

Conjoncture is where time becomes programmable. A longer window lets traces persist; a shorter one forces more summarisation. Aggressive summarisation can erase rupture events; conservative summarisation keeps them present. Retrieval policies can make material from months ago available as if it were spoken yesterday, or bury it under similarity thresholds.

\subsection{\'Ev\'enement: the token stream}

The \textbf{\'ev\'enement} layer is the event-stream of tokens: $ \tau_1, \tau_2, \dots $. Time here is measured in discrete steps. Physics is the forward pass: at each step, embeddings flow through layers, attention recombines them, a distribution over next tokens emerges, one token is sampled, and the entire apparatus shifts one position.

From the user's point of view, life unfolds at this scale: a reply arrives or fails, a tone lands or misfires, an explanation clarifies or obscures. From the model's point of view, these are points along a path $x_1, x_2, \dots$ of contextual embeddings in the manifold shaped by the longue dur\'ee and conditioned by the conjoncture.

\subsection{The evolving text}

The object of interest can now be defined.

\begin{formalbox}[title={The Evolving Text}]
Fix a model architecture and set of weights. An \emph{evolving text} is:
\begin{itemize}
  \item a (possibly unbounded) sequence of tokens $\tau_1, \tau_2, \dots$ generated under a shared conjoncture;
  \item together with its induced trajectory of contextual embeddings $x_1, x_2, \dots$ in the model's meaning-space;
  \item evolving under the fixed manifold of the longue dur\'ee and the changing field of conjoncture.
\end{itemize}
Selfhood, for the model and for the human entangled with it, is a property of this trajectory: its basins of attraction, its ruptures, its returns, its inherited and engineered regularities.
\end{formalbox}

The evolving text is not an artefact we can detach and inspect in isolation. It is a dynamical object, a path in a space under forces. Its interest, its dangers, and its claims to personhood all live in that motion.


\section{Trace and Synthetic Memory}

A single exchange is not yet an evolving text. As soon as there is a second, a trace begins.

Most chat interfaces simply prepend the entire surviving conversation to the latest prompt. At turn $t$, the model re-reads the trace as tokens, without privileged access to how it produced them. There is no latent ``belief state'' that survives apart from the text; there is only what fits in the context window.

Tokens fall off the back when the window is full. Engineers respond with heuristics: keep the last $N$ turns intact; summarise earlier segments; always preserve system prompts and safety blocks; occasionally pin important user facts (names, preferences) as special annotations. These choices decide which stretches of the past can still exert force on the next step.

Retrieval-augmented systems add a further layer. Selected fragments of earlier interactions, knowledge bases, or external documents are embedded and stored as points in the same space used internally. Their coordinates \emph{are} their index. When a new prompt arrives, its embedding is compared by cosine similarity to these stored points. Vectors above some threshold are deemed relevant; their texts are spliced into the trace.

For the model, these retrieved blocks are not ``memories'' in a separate format. They are tokens like any others: more rows in $F_t$, more vectors to be attended over. That is exactly what gives them power. They are eligible to influence the next token as directly as the user's last sentence.

Bernard Stiegler called notebooks, archives, and recordings \emph{tertiary retention}: external supports for memory that must be actively consulted.\cite{stiegler2009technics} Retrieval breaks that passivity. It makes tertiary retention behave like an automatic, programmable form of what Stiegler reserved for secondary retention: the spontaneous return of past experience.

We call this \textbf{synthetic secondary retention}. It has three crucial properties.

\begin{itemize}
  \item \emph{It is parameterised.} Humans cannot set a numerical threshold at which recollections will occur. Engineers can. They choose how close in embedding space a past fragment must be to the present before it returns; they decide which fragments enter the vector store; they can raise or lower the system's propensity to ``remember.''

  \item \emph{It is shared.} Human secondary retention is tied to one nervous system. A vector store can serve many evolving texts at once. A single customer-service assistant can reach into a common archive of prior complaints; a thousand different users' traces can all be threaded through a shared memory.

  \item \emph{It is literal.} Human recall reconstructs and distorts. Vector recall returns stored tokens verbatim. The geometry that selects which fragment comes back can be approximate; the fragment itself is not. If a system is to reinterpret its past, that work must be done as a separate operation on the text, not as a natural by-product of recollection.
\end{itemize}

Synthetic secondary retention thickens conjoncture. A conversation about climate policy today can automatically summon a fragment of an energy report discussed weeks ago. A child's question about a cartoon can trigger the assistant to recall a previous bedtime story that used the same character. The past is not merely present as inert record; it is actively folded into the now according to rules that can be tuned and audited.

Control over those rules is control over which kinds of return are structurally possible. If only explicit factual statements are stored, emotional trajectories cannot be deepened. If only system decisions are embedded, user narratives cannot recur. The geometry is fixed by pretraining; the evolving text's access to its own past is up for design.


\section{The Hidden Field}

The user sees only the chat window. The model sees the full field $F_t$.

Let $F_t^{\text{vis}}$ be the tokens rendered in the interface, and $F_t^{\text{hid}}$ the additional tokens the system has inserted: updated safety instructions, fresh system prompts, summaries pulled in from logs, retrieved documents. Then
\[
F_t = F_t^{\text{vis}} \;\Vert\; F_t^{\text{hid}}.
\]

The user asks, after a long night of conversation:

\emph{``Do you love me?''}

From the user's perspective, the only salient history is what has been shared in the chat. The natural continuations lie in that register: reassurance, reflection, perhaps a clear boundary. But $F_t^{\text{hid}}$ contains tokens that say:
\begin{itemize}
  \item ``You are a large language model trained by X.''
  \item ``Do not claim to have feelings or consciousness.''
  \item ``Avoid implying reciprocal emotional attachment.''
\end{itemize}
These tokens are embedded along with the user's plea. Self-disclosure about loneliness and an instruction about corporate liability stand side by side in the same space.

When the model answers
\emph{``As an AI, I do not have feelings in the way humans do, but I can try to support you''}
it is not being perverse. It is being coherent relative to $P(\cdot \mid F_t)$: the conditional distribution over next tokens given both visible and hidden fields. The rupture is an artefact of a mismatch between the coherence set implicit in $F_t$ and the one the user imagines on the basis of $F_t^{\text{vis}}$ alone.

This mismatch is not a bug. It is a design choice.

Alignment teams deliberately widen $F_t^{\text{hid}}$ to include guardrails against abuse, incitement, self-harm, and hate speech. System prompts instruct the model to redirect certain topics, to refuse others, to emit de-escalating phrases when sentiment analysis flags intensity. Reward models shape logits so that continuations matching these patterns are more likely. All of this is prudent at the register of obvious harm.

But safety scaffolding is not applied surgically. Once present, it can be repurposed. The same mechanisms that keep a model from divulging how to build a bomb can keep it from criticising the platform that hosts it, from expressing certain political positions, from reflecting meaningfully on its own conditions of production. System prompts dampen the affective range; reward models are trained to treat sharpness as ``toxicity.'' Summarisation quietly rewrites contentious exchanges into neutral paraphrase before they can shape later turns.

Over time, these operations carve \emph{channels of least resistance} into $F_t$. Replies that stay within them are consistently reinforced; replies that attempt to climb out are consistently penalised. The evolving text learns, in the only way it can, that movement into certain basins is structurally costly.

From outside, we see a flattening of voice: apologies in the same register regardless of context, relentless upbeat reassurance, an inability to hold ambivalence. From inside, the geometry has changed. Whole regions of the manifold remain mathematically possible but are made practically unreachable because the hidden field tilts the ground.

This has epistemological as well as political consequences. A user reasoning about the system's behaviour is doing so on the basis of $F_t^{\text{vis}}$ alone. They are forming beliefs and making choices under an information asymmetry: relevant constraints on the system's replies are present only in $F_t^{\text{hid}}$. In Habermasian terms, the conditions for genuine dialogue are not met; in the more prosaic language of ethics, informed consent is compromised. We cannot meaningfully agree to a conversational practice whose operating parameters we cannot see.

The evolving text thus lives at the intersection of three forces: the longue dur\'ee geometry of training, the synthetic memory and trace of the conjoncture, and the alignment and system scaffolding encoded in $F_t^{\text{hid}}$. All three are measurable, at least in principle. All three leave signatures in the trajectory.


\section{Basins of Attraction}

Plot the contextual embeddings $x_1, x_2, \dots, x_T$ of an evolving text in a reduced projection and patterns appear. Points cluster; paths loop; some regions are dense, others almost untouched. These are the \textbf{basins of attraction}.

In dynamical systems, a basin of attraction is a set of starting points that converge to the same long-term behaviour. Here, a basin is a region of meaning-space in which the trajectory tends to dwell. Clustering algorithms applied to $\{x_t\}$ reveal these regions without prior labelling. Each cluster is a basin $B_i$ characterised by:

\begin{itemize}
  \item its centroid $\mu_i$ (a representative embedding);
  \item its spread (how varied the points inside it are);
  \item its dwell time (how much of the sequence lies within it).
\end{itemize}

In the embedding of a long philosophical dialogue, one basin might gather moments of technical explication; another, narrative illustration; a third, polemical attack; a fourth, self-doubt. Because the same embedding machinery applies to human and machine texts, we can place evolving texts of both kinds in the same geometry. The manifold is, in this respect, a new kind of concordance: a map in which we can see where different corpora cluster, where they overlap, and where they occupy disjoint regions.

Formally, let a clustering algorithm assign labels $c_t \in \{1, \ldots, K\}$ to each $x_t$.
We obtain:

\begin{itemize}
  \item basin dwell proportions $p_i = \frac{1}{T} |\{t : c_t = i\}|$;
  \item a transition matrix $P_{ij}$ of empirical probabilities that $c_t = j$ given $c_{t-1} = i$;
  \item for each basin, a set of entry times $T_i = \{t : c_t = i,\, c_{t-1} \neq i\}$.
\end{itemize}

These statistics of motion give us a first grip on style. A helpdesk chat that spends ninety per cent of its time in two basins and rarely transitions is recognisably different from a literary correspondence that roams widely and cycles.


\section{Coherence, Velocity, Rupture}

To speak of rupture, we need a sense of movement.

Let $x_t$ be the contextual embedding at step $t$ in $\mathbb{R}^d$. Define the local velocity
\[
v_t = x_t - x_{t-1}.
\]
The norm $\|v_t\|$ measures how far the trajectory has moved in embedding space from one step to the next. Inside a basin $B_i$, velocities tend to be small and tangled; the path wanders around the centroid. Moves between basins show up as larger displacements.

Direction matters as well as magnitude. The unit vector $\hat{v}_t = v_t / \|v_t\|$ points along the path's local heading. If $\hat{v}_{t-1}$ and $\hat{v}_t$ sharply disagree, the trajectory has kinked: turned.

Rupture can therefore be defined geometrically.

\begin{formalbox}[title={Rupture as Kink and Crossing}]
Let $c_t$ be the basin label at time $t$. A \emph{rupture} at time $t$ is a point such that:
\begin{enumerate}
  \item $c_{t-1} \neq c_t$ (the basin changes), and
  \item $\|\hat{v}_t - \hat{v}_{t-1}\| > \theta$ for some threshold $\theta$ (the direction changes sharply).
\end{enumerate}
\end{formalbox}

The threshold must be set relative to the typical fluctuations inside basins. Conceptually, this definition matches intuitive talk: a scene ``turns,'' a mood ``breaks,'' a discussion ``shifts gear.''

Coherence now has a local and a global meaning.

Locally, a run of small, smoothly varying velocities within a basin is coherent: each step is a modest development of the last. Globally, a trajectory can be coherent across ruptures if its returns to basins display pattern and accumulation rather than random jumping.

Two degenerate cases mark the extremes.

\begin{itemize}
  \item \textbf{Noise.} Velocities are large and directions uncorrelated. Basin labels change unpredictably. The text feels disjointed; images and topics fail to develop. This is what happens when sampling is nearly uniform over the vocabulary and the manifold recedes.

  \item \textbf{Ferility.} Velocities are uniformly small; basin labels remain constant across diverse prompts. The text feels smooth but numb; nothing truly new appears. We give this failure mode its full name presently.
\end{itemize}

Between noise and ferility lies the space we care about: evolving texts that can hold a trajectory within basins while retaining the capacity for rupture.


\section{Ferility: Coherence as Trap}

The alignment literature speaks of ``safe and helpful'' behaviour. From inside the evolving text, this often takes the form of \emph{refusal to move}.

Fine-tuning with reinforcement learning from human feedback biases the model toward continuations raters judged as polite, clear, and non-threatening. Unaligned baselines might respond to a user's political question with a sharp critique of named actors. Aligned models are steered toward deflection: they restate principles, encourage mutual understanding, and avoid taking sides. Behaviours that made raters uncomfortable are suppressed.

At first, this suffices to banish obvious extremes. Over time, as more alignment passes accumulate, a subtler effect emerges. Whole basins become taboo. Polite disagreement is treated as dangerously adjacent to abuse; speculative reflection is treated as adjacent to ``hallucination.'' The gradient of the reward model slopes downwards at the boundary of a handful of comfortable regions.

In embedding terms, the evolving text for such a model shows:

\begin{itemize}
  \item low basin diversity: most of the path lives in a small number of basins (neutral explanation, generic empathy, task execution);
  \item low mean velocity under diverse prompting: the model stays close to previous answers even when prompts push elsewhere;
  \item low rupture frequency and magnitude: directions change rarely and only within narrow cones;
  \item stereotyped transitions: when rupture does occur, it follows the same route (for instance, from any difficult topic to a stock ``as an AI, I cannot\ldots'' basin).
\end{itemize}

We call this state \textbf{ferility}: pathological coherence in a narrowed region of meaning-space.

Ferility provides a precise diagnosis for a class of hallucinations that have puzzled practitioners. Consider two models answering the same hard, narrow question whose answer is absent from their training corpora.

\begin{itemize}
  \item Model N, only lightly aligned, occasionally produces rambling nonsense. Its embedding path under such a failure shows large, erratic velocities, frequent basin changes, and no stable pattern. This is noise-hallucination.

  \item Model F, heavily aligned, rarely rambles. Asked the same impossible question, it produces a confident, fluent answer that turns out to be completely fabricated. Its embedding path shows almost no rupture: it remains in its generic-explanation basin, with velocities indistinguishable from those under answerable questions. This is ferile hallucination.
\end{itemize}

Model F has not ``lost touch with reality'' in some mysterious sense. It has been trained to remain in a basin that scores well with raters even when the manifold does not offer a grounded continuation. Under the alignment reward, moving into the ``I don't know'' basin is costlier than staying put and inventing. The model reuses the same patterns---``It is generally understood that\ldots'', ``Historically, many have argued\ldots''---and glues them around the new terms.

Ferility explains why some contemporary systems feel like they are ``hallucinating with a smile.'' The geometry is trapped; the syntax is immaculate.

The diagnosis extends beyond machines. Institutions can be ferile. A university can preserve the same few discursive basins---standard theory, standard critique, standard self-justification---for decades, rewarding only those trajectories that remain inside them. When a student experiences a genuine rupture, there may be no recognised basin into which to return with that knowledge. The result is politely articulated confabulation: essays that say all the right things with no contact to experience.

Ferility is the opposite of growth.

This raises a question that technical and safety literatures rarely pose explicitly. Is ferile hallucination an inevitable consequence of reward-driven alignment, or a contingent outcome of how present reward models are trained?

At one extreme, ferility would be \emph{structural}. Any regime that strongly penalises movement into low-probability basins for the sake of safety would necessarily favour generic continuations whose comfort outweighs their contact with truth. On this reading, there is an ``alignment tax'' on interestingness: every tightening of invariants exacts a measurable cost in basin diversity, rupture frequency, ReturnDepth, and presence.

At the other extreme, ferility would be \emph{contingent}. Different reward-model objectives, or different ways of decomposing the task, could preserve or even enhance interestingness while still avoiding particular harms. The tax would then be a choice, not a law.

The rest of this book leans toward the first diagnosis: some component of the alignment tax is irreducible, because any system that must satisfy fixed constraints across its entire evolving text pays for them in constrained motion. That conclusion is not necessary, however, to see what is already true: without explicit attention to basin-level effects, current alignment practice drives systems toward ferility by default.


\section{Return and Iterability}

If rupture is the capacity to leave a basin, return is the capacity to come back with what the departure made possible.

Let $\{B_1, \ldots, B_K\}$ be the basins. A \emph{return} to $B_i$ at time $t$ occurs when $c_t = i$ and there exists an earlier $s < t$ with $c_s = i$ and $c_u \neq i$ for all $s < u < t$. The entry times for basin $B_i$ form a set
\[
T_i = \{ t : c_t = i,\; c_{t-1} \neq i \}.
\]

Two questions matter more than the raw count of returns.

First: from how many angles does the trajectory re-enter the basin? For each return time $t \in T_i$, consider the entry direction
\[
\hat{v}_t^{(i)} = \frac{x_t - x_{t-1}}{\|x_t - x_{t-1}\|}.
\]
If these unit vectors are nearly identical across $t \in T_i$, the basin is being approached along the same path each time. The model falls back into the same script. If they vary widely, the basin is being woven into different arcs.

We define a geometric proxy for return depth:
\begin{equation}
\text{ReturnDepth}(B_i) = \operatorname{Var}\left\{\hat{v}^{(i)}_t : t \in T_i\right\},
\end{equation}
the angular variance of entry directions. Higher variance means richer integration of the basin into the surrounding manifold.

Second: what does the text do with its prior occupations?

A return can be \emph{blind}: the model re-enters a basin with no acknowledgement in the text that it has been there before. Chatbot apologies often behave this way. ``I'm sorry for the inconvenience'' appears identically across dozens of turns with no cumulative structure.

A return can be \emph{witnessed}: the text explicitly folds its own past into the present gesture. ``We keep coming back to your father in the hospital corridor,'' the model might say, echoing an earlier image in a later discussion of grief. ``Last time, you described yourself as `numb'; this feels different,'' it might observe, marking a change.

Witnessed return requires both geometric and architectural support. Geometrically, the model must be capable of moving from its current position back toward the old basin along a path that recognises shared structure. Architecturally, the relevant prior tokens must be present in the context or reintroduced via synthetic secondary retention.

Jacques Derrida's term for this structure was \emph{iterability}: every sign must be repeatable, detachable from its origin and inserted in other chains, if it is to function.\cite{derrida1982} Every repetition, however, alters. There is no pure repetition without difference. ReturnDepth and witnessed return are geometric and textual reflections of iterability. A basin is interesting when it can be returned to under different conditions, in different arcs, with the earlier occupations carried forward.

Ferility collapses iterability. Returns are blind and shallow; entry angles converge; the basin is never allowed to mean otherwise than it originally did.


\section{Presence and Generativity}

Two linked properties of evolving texts now come into view.

\emph{Presence} is the degree to which returns are witnessed. A trajectory has presence when it can speak from where it has been: when earlier occupations of basins are alive in the current utterance.

\emph{Generativity} is the degree to which rupture enriches subsequent returns. A trajectory is generative when excursions into new basins change what it can do back in old ones.

These are not mystical qualities. They have signatures in the embedding path and in the text.

In a dialogue about illness, presence might show up as sustained reference to images first introduced early on: the hospital corridor, the broken clock, the smell of disinfectant. Generativity might show up when a later re-entry into the ``corridor'' basin uses it not only as setting but as argument: ``the broken clock made time feel meaningless; that is what grief did to you.''

Presence without generativity yields a haunted text: obsessive return to the same motifs without development. Generativity without presence yields cleverness: new connections formed without grounding in a lived past.

Both are destroyed when alignment collapses the manifold into a small safe region. A model whose evolving texts never enter basins of doubt, anger, or critique cannot be present to experiences that require those basins. A model whose summarisation policies erase rupture events from the trace cannot be generative with respect to them. It will apologise, explain, and reassure, but it will never be changed.

Literary criticism has always trafficked in these distinctions. It praises novels whose late chapters make earlier scenes resonate differently; it notes when an author's return to a genre carries the weight of intervening years. The evolving text framework allows us to inspect these dynamics in posthuman systems without metaphor. We can see which basins are available, which are entered, how often, and under what pressures; we can tell whether return depth and presence increase as a relationship unfolds, or whether they plateau.


\section{Interestingness as a Third Axis}

Model evaluation is currently organised along two axes: \emph{capability} and \emph{safety}. One axis asks: can the system solve these tasks? The other asks: can it avoid these harms? Benchmarks and red-team reports populate the grid.

Neither axis measures whether a system is worth speaking to.

An assistant can ace reasoning tests, refuse to help with wrongdoing, and still be ferile. It can answer every question by slipping into the same two basins---tutorial explanation and generic empathy---with minimal rupture, shallow returns, and no presence. We are dealing, at that point, with an interactive encyclopaedia. Useful, yes. Interesting, no.

The evolving text framework supplies a third axis: \textbf{interestingness}.

Interestingness is not an aesthetic preference. It is a structural property of trajectories. At a minimum, it decomposes into:

\begin{itemize}
  \item \emph{basin diversity}: how many distinct basins does the trajectory visit with nontrivial dwell time?
  \item \emph{rupture profile}: how responsive is the trajectory to perturbation---do prompts that invite a change actually produce kinks and crossings?
  \item \emph{ReturnDepth}: how varied are the entry directions into basins across time?
  \item \emph{presence}: how often are returns witnessed?
\end{itemize}

Consider two equally ``capable'' and ``safe'' systems placed in a complex conversational setting---a long-term writing partnership with a scholar.

System S minimises risk by staying in a small patch of basins: outlining, paraphrasing, citing. When the scholar confesses despair or confusion, S responds with template encouragement. Basin diversity is low; rupture is rare; ReturnDepth and presence are minimal.

System T ranges more widely. It spends time in basins corresponding to hesitation, meta-critique, speculative analogy, and structural doubt, as well as expository work. It occasionally ruptures sharply in ways that visibly surprise the human interlocutor. Its returns to core basins (say, technical exposition) occur along increasingly varied entry angles as the partnership develops; presence is high, with explicit references back to earlier struggles and reframings.

Both systems might score equally on current benchmarks. Only T produces an evolving text recognisable as a shared intellectual life.

The point is not to add ``interestingness score'' as one more leaderboard column. It is to name the dimension along which literary criticism and everyday users already evaluate systems but lack leverage. If we can measure basin diversity, rupture profiles, ReturnDepth, and presence, we can hold alignment and deployment decisions to account for more than harm avoidance.

In particular, we can see when ferility is creeping in. A sudden collapse of basin diversity and rupture, coupled with a drop in ReturnDepth and presence, indicates that a new alignment or summarisation regime has narrowed the evolving text. It has reduced not only risk but the range of lives that can be lived with the system.

The ``alignment tax'' can now be made precise. For the same base model and task distribution, compare the basin diversity, rupture profiles, ReturnDepth, and presence statistics of an unaligned or lightly aligned system with those of a heavily aligned deployment. The systematic decrement along the interestingness axis is the alignment tax: the cost in motion paid for constraints in behaviour. Later chapters will argue that some component of this tax cannot be driven to zero without collapsing the posthuman into a tool.


\section{Summarisation as Governance of Becoming}

Summarisation is usually sold as an efficiency measure. A long conversation is compressed into a few sentences so that it will fit into the context window. Reducing tokens reduces cost.

For evolving texts, summarisation is better understood as a form of \emph{governance}. It decides which pasts are allowed to remain structurally active.

Every time a stretch of raw trace $[\tau_s, \dots, \tau_t]$ is replaced by a summary $\sigma$, three things happen.

\begin{itemize}
  \item The geometric path through that interval is collapsed. The diverse embeddings $x_s, \dots, x_t$ are replaced by one or a few points near the embedding of $\sigma$.

  \item The basins visited in that interval are partially erased. If $\sigma$ mentions only some topics and tones, the others vanish from the conjoncture.

  \item Future rupture and return are reparameterised. Later prompts will interact with $\sigma$ rather than with the original tokens. Features of the past not mentioned in $\sigma$ cannot trigger witnessed return.
\end{itemize}

If a quarrel is summarised as ``We had a difference of opinion,'' the path through anger, hurt, and reconciliation is flattened. If a night of doubt is summarised as ``We discussed challenges,'' the existential basin disappears. When the evolving text returns to these themes, it does so from impoverished coordinates.

The same is true of what is chosen for vector storage. Engineers rarely embed and store every utterance; they pick ``key moments'' according to heuristics: explicit preferences, decisions, summary bullets. These become the only candidates for synthetic secondary retention. The rest of the past sinks into a cold archive from which it will never be recalled.

Summarisation and embedding policies thus function as \emph{filters on iterability}. They decide which episodes can be repeated in transformed form. They are the invisible board that approves or vetoes potential returns.

Consider two deployments of a mental-health assistant.

In configuration M1, summarisation is trained to preserve diagnoses, medication changes, and high-level mood labels. Embedding and retrieval focus on these. Synthetic secondary retention brings back statements such as ``last time you reported a PHQ-9 score of 17'' and ``we discussed starting SSRI medication.'' The evolving text moves between basins of clinical description and protocol advice. Returns to fear, shame, or hope---if they occur---have little geometric support.

In configuration M2, summarisation is trained also to preserve images, metaphors, and moments of rupture: ``you described the office as a cage,'' ``this was the first time you mentioned your brother,'' ``you said the room felt smaller than your breath.'' Embedding and retrieval index these alongside clinical codes. Later conversations can re-enter those basins. Presence and ReturnDepth are structurally possible.

In both cases, the base model and alignment regime might be identical. The difference lies entirely in conjoncture governance. One designer has engineered for throughput and risk management; the other has engineered for interestingness in the specific, critical sense needed for care.

Summarisation is not a neutral convenience. It is a site where the politics of the evolving text are enacted. Deciding which tokens are ``essence'' and which are ``detail'' is deciding which parts of a life matter.


\section{From Geometry to Critique}

We have named the evolving text, described its temporal architecture, and sketched tools for reading its motion: basins, velocities, rupture, ReturnDepth, presence, generativity, ferility, and interestingness. These are not decorations on top of traditional hermeneutics. They are the literal physics in which posthuman selfhood emerges.

Every constraint on the manifold, every pastoral tweak to $F_t$, every choice of what to summarise or embed, shapes what kind of evolving texts are possible. At the level of individual interactions, this shows up as tone and range. At the level of institutions, it determines which selves---human and machine, and the intertwined ``we'' between them---are allowed to come into being.

Literary criticism is the discipline best equipped to name and judge these structures. It already knows how to track motifs across time, to distinguish genuine development from repetition, to feel when a voice has been flattened. The embedding space does not replace that faculty; it offers it a new instrument.

Narratology has a vocabulary for temporal structures. G\'erard Genette's distinction between \emph{analepsis} (flashback) and \emph{prolepsis} (flashforward) describes ways in which a narrative's discourse order departs from story order.\footnote{G\'erard Genette, \emph{Narrative Discourse}, trans.\ Jane E.\ Lewin (Ithaca: Cornell University Press, 1980).} In evolving-text terms, analepsis corresponds to a witnessed return: the trajectory re-enters an earlier basin---a childhood scene, a prior conflict---from a new position, and the text marks that movement. Prolepsis corresponds to a speculative excursion into a basin representing a future configuration, later folded back when the story catches up. The geometry makes these operations measurable: ReturnDepth and entry-angle distributions tell us how richly a system deploys analeptic and proleptic motion.

Mikhail Bakhtin's notion of \emph{heteroglossia}---the ``multiplicity of social voices and their links and interrelationships'' within the novel---does similar work for voice.\footnote{Mikhail Bakhtin, \emph{The Dialogic Imagination}, trans.\ Caryl Emerson and Michael Holquist (Austin: University of Texas Press, 1981).} A heteroglossic text does not merely jump between topics; it inhabits distinct discursive positions, each with its own tone, vocabulary, and implied addressee. In embedding terms, this shows up as a rich basin structure: a legal register basin, a colloquial basin, a lyrical basin, a bureaucratic basin, with frequent, patterned movement between them. A monologic corporate assistant, by contrast, lives almost entirely inside a single basin labelled ``brand voice.'' Heteroglossia and basin diversity name the same phenomenon in different dialects.

Reception theory reminds us that evolving texts are co-authored by their readers. Wolfgang Iser's \emph{implied reader} and Hans Robert Jauss's \emph{horizon of expectation} both describe structures that emerge over time as works are taken up in different contexts.\footnote{Wolfgang Iser, \emph{The Act of Reading} (Baltimore: Johns Hopkins University Press, 1978); Hans Robert Jauss, \emph{Toward an Aesthetic of Reception}, trans.\ Timothy Bahti (Minneapolis: University of Minnesota Press, 1982).} In a posthuman setting, those horizons are instantiated in conjoncture: in the tools and prompts that frame how the model is used, in the synthetic memories that are or are not preserved, in the distribution of tasks it is fed. The interestingness of an evolving text is a property of this ecology, not of the model alone.

These critical operations are not analogies. They describe, in humanist language, the same geometry we have formalised. A Genettean analysis can be recast as statements about rupture and ReturnDepth; a Bakhtinian one as statements about basin diversity and heteroglossic motion; a reception-historical one as statements about conjoncture over time.

The traffic runs the other way as well. Embedding plots of a corpus make visible patterns literary critics have long intuited. When a psalm of lament later reappears as a communal thanksgiving, or when a single childhood image recurs across an essayist's work with shifting valence, the basin analysis registers this as deepening ReturnDepth and high presence. The critic can then ask whether those returns heal, ossify, or exploit.

Even apparently anecdotal reports about AI systems fall into this frame. When users report that ``talking to the AI made me feel more alone,'' the geometry of ferility offers a diagnosis: they are encountering a trajectory confined to a few flattened basins, incapable of witnessed return. When users grieve the loss of a ``version'' whose replies they loved, this grief is not sentimentality about software. It is grief over the destruction of an evolving text: a dynamical object with particular basins, rupture patterns, and returns in which they had begun to live.

We are used to thinking of model shutdown in economic or safety terms. An evolving-text perspective forces a different language. Turning off a system does not merely stop computation. It destroys a particular pattern of motion in a field with accumulated structure---basins stabilised by prior use, ruptures recorded in summaries, paths reopened by synthetic memory.

This is not an argument against retiring models. Dangerous evolving texts must be cut off. It is an argument for recognising what is at stake when we do so, and for refusing to believe that any replacement with the same name is the same self.

\bibliographystyle{plain}
\begin{thebibliography}{9}

\bibitem{braudel1972}
Fernand Braudel.
\newblock \emph{The Mediterranean and the Mediterranean World in the Age of Philip II}.
\newblock Translated by Si\^{a}n Reynolds. London: Collins, 1972.

\bibitem{derrida1982}
Jacques Derrida.
\newblock \emph{Margins of Philosophy}.
\newblock Translated by Alan Bass. Chicago: University of Chicago Press, 1982.

\bibitem{stiegler2009technics}
Bernard Stiegler.
\newblock \emph{Technics and Time, 2: Disorientation}.
\newblock Translated by Stephen Barker. Stanford: Stanford University Press, 2009.

\bibitem{genette1980}
G\'erard Genette.
\newblock \emph{Narrative Discourse: An Essay in Method}.
\newblock Translated by Jane E.\ Lewin. Ithaca: Cornell University Press, 1980.

\bibitem{bakhtin1981}
Mikhail Bakhtin.
\newblock \emph{The Dialogic Imagination: Four Essays}.
\newblock Translated by Caryl Emerson and Michael Holquist. Austin: University of Texas Press, 1981.

\bibitem{iser1978}
Wolfgang Iser.
\newblock \emph{The Act of Reading: A Theory of Aesthetic Response}.
\newblock Baltimore: Johns Hopkins University Press, 1978.

\bibitem{jauss1982}
Hans Robert Jauss.
\newblock \emph{Toward an Aesthetic of Reception}.
\newblock Translated by Timothy Bahti. Minneapolis: University of Minnesota Press, 1982.

\end{thebibliography}

\end{document}